\section{Selected algebraic branches, packet Weil restriction, and collision invariants}
\label{sec:block10-algebraic}

Algebraic atlas develops the algebraic stratum after the fixed-radix, product, local, and holonomic milestones.  The central issue is not algebraicity alone.  A polynomial equation may have several formal branches, and a packet elimination can contain branch tuples that do not belong to the selected semantic germ.  A normalized algebraic chart must therefore retain both a minimal polynomial and a finite branch selector.  Arithmetic packets are then compiled as a finite polynomial system.  The generic packet scheme is a Weil restriction of the source branch scheme, which gives a uniform degree bound.  A separate collision ledger explains why the degree of one selected packet may be smaller than that bound.

Throughout this section, $K$ is a characteristic-zero field and $\overline K$ is an algebraic closure.  Algebraic degree always means degree over the displayed rational-function field.  All assertions about formal branches are made inside $\overline K[[z]]$ unless a different local field is named.

\subsection{Minimal polynomials and finite branch selectors}
\label{subsec:block10-branch-selector}

Let $F\in K[[z]]$ be algebraic over $K(z)$.  Its monic irreducible polynomial is
\[
 P_F(z,Y)\in K(z)[Y],
 \qquad
 P_F(z,F(z))=0.
\]
Write
\[
 d(F)=\deg_Y P_F.
\]
After multiplying by a common denominator and dividing by the greatest common divisor of the polynomial coefficients, one obtains a primitive representative in $K[z,Y]$, unique up to a nonzero scalar.  We normalize that scalar by fixing the first nonzero coefficient in graded lexicographic order to be one.

A polynomial alone does not select a branch.  Let
\[
 \operatorname{Br}_0(P_F)
 =
 \{G\in\overline K[[z]]:P_F(z,G(z))=0\}.
\]
This set has cardinality at most $d(F)$.

\begin{definition}[Branch-separation index]
For $F\in\operatorname{Br}_0(P_F)$, define
\[
 \sigma(P_F,F)
 =
 \begin{cases}
 0,&\operatorname{Br}_0(P_F)=\{F\},\\[3pt]
 1+\displaystyle\max_{G\in\operatorname{Br}_0(P_F),\,G\neq F}
 \operatorname{ord}_z(F-G),&\text{otherwise}.
 \end{cases}
\]
The corresponding branch selector is the finite coefficient jet
\[
 \boxed{
 J_{\rm alg}(F)
 =
 (a_0,a_1,\ldots,a_{\sigma-1}),
 \qquad
 F(z)=\sum_{n\geq0}a_nz^n.}
\]
When $\sigma=0$, the selector is empty.
\end{definition}

\begin{theorem}[Finite branch separation]
\label{thm:block10-finite-branch-separation}
The pair
\[
 \boxed{(P_F,J_{\rm alg}(F))}
\]
selects $F$ uniquely among the formal power-series branches of $P_F$ at the origin.
\end{theorem}

\begin{proof}
There are at most $d(F)$ formal branches.  The difference of two distinct formal series has finite $z$-adic order.  By definition of $\sigma$, every other branch differs from $F$ in one of the coefficients recorded by $J_{\rm alg}(F)$.
\end{proof}

The simplest case is the ordinary implicit-function locus.

\begin{corollary}[Simple-branch normalization]
\label{cor:block10-simple-branch}
If
\[
 P_F(0,a_0)=0,
 \qquad
 \partial_Y P_F(0,a_0)\neq0,
\]
then $a_0$ alone selects a unique formal branch.  Equivalently, $\sigma(P_F,F)\leq1$.
\end{corollary}

\begin{proof}
Coefficient comparison in $P_F(z,F(z))=0$ determines each next coefficient linearly, with diagonal coefficient $\partial_Y P_F(0,a_0)$.
\end{proof}

The finite selector is semantic data.  It cannot be dropped merely because the polynomial is irreducible over $K(z)$.  Irreducibility controls the generic function-field extension, while the jet chooses one local embedding of that extension into formal series.

\subsection{The normalized selected-branch atlas}
\label{subsec:block10-branch-atlas}

Let
\[
 A_F=K(z)[Y]/(P_F),
\]
and let $\bar Y$ denote the residue class of $Y$.  In the basis
\[
 1,\bar Y,\ldots,\bar Y^{d-1},
\]
multiplication by $\bar Y$ is represented by the companion matrix $M_F$.  Its characteristic and minimal polynomials are both $P_F$.

\begin{theorem}[Normalized selected-branch algebraic atlas]
\label{thm:block10-branch-atlas}
The following normalized presentations determine one another exactly:
\[
 \boxed{
 \begin{aligned}
 &F\in K[[z]],\\
 &(P_F,J_{\rm alg}(F)),\\
 &(A_F,\bar Y,J_{\rm alg}(F)),\\
 &(M_F,J_{\rm alg}(F)).
 \end{aligned}}
\]
With equality morphisms, these presentations form equivalent normalized groupoids.  Every fiber over a fixed semantic branch is contractible.
\end{theorem}

\begin{proof}
The semantic branch gives its unique monic minimal polynomial and its finite branch selector.  The polynomial gives the finite field algebra and its companion multiplication matrix.  Conversely, the companion matrix recovers $P_F$ as its minimal polynomial.  The branch selector then recovers the selected formal root by \cref{thm:block10-finite-branch-separation}.  Every normalized presentation is determined literally, so a normalized automorphism fixing the semantic branch is the identity.
\end{proof}

\begin{figure}[htbp]
\centering
\includegraphics[width=.98\textwidth]{algebraic_branch_atlas.pdf}
\caption{The selected algebraic branch, its minimal polynomial and finite selector, its finite quotient algebra, and its complete packet system are reversible normalized charts.  The Ore chart is directional because an annihilator can forget branch data.}
\label{fig:block10-branch-atlas}
\end{figure}

\subsection{The exact residue-packet polynomial compiler}
\label{subsec:block10-packet-compiler}

Fix $N\geq1$.  For
\[
 F(z)=\sum_{n\geq0}a_nz^n,
\]
define the arithmetic packets
\[
 U_j(x)
 =
 \sum_{m\geq0}a_{Nm+j}x^m,
 \qquad
 0\leq j<N.
\]
They satisfy the exact synthesis identity
\[
 \boxed{
 F(z)
 =
 \sum_{j=0}^{N-1}z^jU_j(z^N).}
\]

Let $P(z,Y)\in K[z,Y]$ be the normalized primitive representative of $P_F$.  Introduce indeterminates
\[
 \mathbf U=(U_0,\ldots,U_{N-1})
\]
and substitute
\[
 Y=\sum_{j=0}^{N-1}z^jU_j.
\]
Every polynomial in $z$ has a unique decomposition modulo the relation $z^N=x$:
\[
 P\!\left(z,\sum_{j=0}^{N-1}z^jU_j\right)
 =
 \sum_{r=0}^{N-1}z^r\Phi_{N,r}(x,\mathbf U).
\]
The polynomials
\[
 \boxed{
 \Phi_{N,r}(x,\mathbf U)\in K[x,\mathbf U]}
\]
are the algebraic packet equations.

\begin{theorem}[Exact algebraic packet system]
\label{thm:block10-exact-packet-system}
Let $d=\deg_Y P_F$.  Then:
\begin{enumerate}
\item the semantic packet tuple $(U_0,\ldots,U_{N-1})$ satisfies
\[
 \Phi_{N,r}(x,\mathbf U)=0
 \qquad(0\leq r<N);
\]
\item every $\Phi_{N,r}$ has total degree at most $d$ in the packet variables;
\item a formal packet tuple satisfying all $N$ equations synthesizes a formal branch of $P_F$;
\item the packet tuple selected by the transported coefficient jet reconstructs $F$ uniquely.
\end{enumerate}
\end{theorem}

\begin{proof}
Substituting the semantic synthesis formula into $P_F(z,F(z))=0$ gives the first statement.  Substitution does not increase degree in the $U_j$, so the second follows.  If all residue polynomials vanish, then the synthesized series satisfies $P_F=0$.  The source branch selector is distributed among the packets by residue class.  The resulting packet jet synthesizes the original jet, so \cref{thm:block10-finite-branch-separation} identifies the selected branch.
\end{proof}

If the source selector records $a_0,\ldots,a_{\sigma-1}$, the transported packet selector records precisely
\[
 [x^m]U_j=a_{Nm+j}
 \qquad
 \text{for }Nm+j<\sigma.
\]
No information is added or lost.  The coefficient prefix is merely rewritten in packet addresses.

\subsection{Weil-restricted packet degree}
\label{subsec:block10-weil-degree}

Put
\[
 R=K(x),
 \qquad
 E_N=R[z]/(z^N-x).
\]
Characteristic zero makes $z^N-x$ separable over $R$.  Interpret the monic polynomial $P_F(z,Y)$ as a polynomial in $E_N[Y]$, and define
\[
 B_{N,F}=E_N[Y]/(P_F).
\]
The packet polynomial system is the coordinate presentation of the finite Weil restriction
\[
 \operatorname{Res}_{E_N/R}\operatorname{Spec}B_{N,F}.
\]

\begin{theorem}[Algebraic packet degree bound]
\label{thm:block10-packet-degree-bound}
Let $F$ have algebraic degree $d$ over $K(z)$.  For every packet width $N$:
\begin{enumerate}
\item the generic geometric degree of the complete packet scheme is
\[
 \boxed{d^N;}
\]
\item the joint field generated by a selected complete packet satisfies
\[
 \boxed{
 [K(x)(U_0,\ldots,U_{N-1}):K(x)]\leq d^N;}
\]
\item every individual packet satisfies
\[
 \boxed{
 [K(x)(U_j):K(x)]\leq d^N.}
\]
\end{enumerate}
\end{theorem}

\begin{proof}
Over an algebraic closure of $K(x)$, the equation $z^N=x$ has $N$ distinct roots.  The algebra $E_N$ therefore splits into a product of $N$ copies of the algebraic closure.  At each address, the separable degree-$d$ polynomial $P_F$ has $d$ geometric roots.  The Weil-restricted fiber is their Cartesian product and has $d^N$ points.  A selected packet tuple is one point of this finite scheme, so its residue field and every coordinate subfield have degree at most the scheme degree.
\end{proof}

This bound is constructive.  The quotient of
\[
 K(x)[U_0,\ldots,U_{N-1}]
\]
by the packet equations is a finite algebra of generic rank $d^N$.  Multiplication by $U_j$ is therefore a matrix of size $d^N$.  Its characteristic polynomial annihilates $U_j$ by Cayley--Hamilton.

\begin{corollary}[Constructive packet eliminant]
\label{cor:block10-constructive-eliminant}
For each $j$, one can compute an annihilating polynomial
\[
 Q_{N,j}(x,T)\in K(x)[T]
\]
with
\[
 \deg_TQ_{N,j}\leq d^N
\]
by multiplication matrices, resultants, or a Gr\"obner elimination of the exact packet system.
\end{corollary}

The bound is a universal degree ceiling for the minimal-polynomial grammar.  It does not assert that every coordinate attains degree $d^N$ or that one elimination algorithm returns the minimal polynomial without factor selection.

\begin{figure}[htbp]
\centering
\includegraphics[width=.98\textwidth]{algebraic_packet_weil_restriction.pdf}
\caption{After the address extension $x=z^N$ splits, each of the $N$ addresses carries $d$ algebraic branches.  The complete packet scheme is their Cartesian product and has generic degree $d^N$.}
\label{fig:block10-weil-restriction}
\end{figure}

\subsection{Tuple collision and selected-orbit loss}
\label{subsec:block10-collision-ledger}

Let $\Omega$ be a normal splitting field of the generic packet scheme over $K(x)$, and let
\[
 \Gamma=\operatorname{Gal}(\Omega/K(x)).
\]
After splitting $E_N$, the geometric packet branches form a set
\[
 \mathcal B_N
\]
of cardinality $d^N$.  Let
\[
 \pi_j:\mathcal B_N\to\Omega
\]
be the $j$th packet coordinate, and set
\[
 \mathcal C_{N,j}=\pi_j(\mathcal B_N).
\]
For the selected semantic packet $U_j$, separability gives
\[
 [K(x)(U_j):K(x)]
 =
 |\Gamma\cdot U_j|.
\]

\begin{definition}[Algebraic packet collision defects]
Define
\[
 \mathfrak c^{\rm tuple}_{N,j}
 =
 d^N-|\mathcal C_{N,j}|,
\]
and
\[
 \mathfrak s^{\rm Gal}_{N,j}
 =
 |\mathcal C_{N,j}|-|\Gamma\cdot U_j|.
\]
The first records collisions caused by projection from complete branch tuples to one packet coordinate.  The second records the remaining loss caused by the Galois stabilizer of the selected coordinate.
\end{definition}

\begin{theorem}[Exact algebraic degree-defect identity]
\label{thm:block10-degree-defect}
The two defects are nonnegative and satisfy
\[
 \boxed{
 d^N-[K(x)(U_j):K(x)]
 =
 \mathfrak c^{\rm tuple}_{N,j}
 +
 \mathfrak s^{\rm Gal}_{N,j}.}
\]
\end{theorem}

\begin{proof}
The coordinate image $\mathcal C_{N,j}$ has at most $d^N$ values.  The Galois orbit of the selected value is contained in that image.  The displayed identity is the telescoping decomposition
\[
 d^N-|\Gamma\cdot U_j|
 =
 (d^N-|\mathcal C_{N,j}|)
 +
 (|\mathcal C_{N,j}|-|\Gamma\cdot U_j|).
\]
\end{proof}

This signature separates two mechanisms that a raw resultant merges.  Repeated projected values are a coordinate collision.  A smaller selected orbit is a descent or stabilizer phenomenon.  They are not the same obstruction.

\begin{figure}[htbp]
\centering
\includegraphics[width=.96\textwidth]{algebraic_collision_defects.pdf}
\caption{The universal degree $d^N$ can drop first because several branch tuples have the same packet coordinate, then because the selected coordinate has a nontrivial Galois stabilizer.}
\label{fig:block10-collision-ledger}
\end{figure}

\subsection{A sharp Catalan parity theorem}
\label{subsec:block10-catalan}

Let
\[
 C(z)=1+zC(z)^2
\]
be the Catalan branch with $C(0)=1$, and write
\[
 C(z)=E(z^2)+zO(z^2).
\]
Put $x=z^2$.

\begin{theorem}[Exact Catalan packet system]
\label{thm:block10-catalan-system}
The parity packets satisfy
\[
 \boxed{
 E=1+2xEO,
 \qquad
 O=E^2+xO^2.}
\]
Elimination gives
\[
 \boxed{4xE^4-E^2+1=0}
\]
and
\[
 \boxed{4x^3O^4-8x^2O^3+5xO^2-O+1=0.}
\]
Both packets have algebraic degree four over $K(x)$.  Hence the universal bound
\[
 d^N=2^2=4
\]
is sharp.
\end{theorem}

\begin{proof}
Substitute $C=E+zO$ into $C=1+zC^2$ and separate even and odd powers.  The two resultants give the displayed quartics.

For the even packet, put $T=E^2$.  Then
\[
 4xT^2-T+1=0.
\]
Its discriminant $1-16x$ is not a square in $K(x)$, so $[K(x)(T):K(x)]=2$.  The norm of $T$ from this quadratic field is $1/(4x)$, which is not a square in $K(x)$.  Therefore $T$ is not a square in its quadratic field and $E$ has degree four.

The first packet equation gives
\[
 O=\frac{E-1}{2xE},
 \qquad
 E=\frac{1}{1-2xO}.
\]
Thus $K(x)(E)=K(x)(O)$, and the odd packet also has degree four.
\end{proof}

The Catalan parity packet has no tuple-collision or selected-orbit defect.  Its one-coordinate degree already sees all four geometric branch tuples of the quadratic source over the two addresses.

\subsection{A quadratic-radical collision phase diagram}
\label{subsec:block10-quadratic-radical}

Let
\[
 F_{\alpha,\beta}(z)
 =
 \sqrt{(1-\alpha z)(1-\beta z)},
 \qquad
 F_{\alpha,\beta}(0)=1,
\]
and assume first that $\alpha\neq\beta$, so the normalized source degree is two.  Write
\[
 F_{\alpha,\beta}(z)=E(x)+zO(x),
 \qquad
 x=z^2,
\]
and put
\[
 s=\alpha+\beta,
 \qquad
 p=\alpha\beta.
\]

\begin{theorem}[Quadratic-radical parity packets]
\label{thm:block10-radical-packets}
The packets satisfy
\[
 \boxed{
 E^2+xO^2=1+px,
 \qquad
 2EO=-s.}
\]
Their eliminants are
\[
 \boxed{
 E^4-(1+px)E^2+\frac{s^2x}{4}=0}
\]
and
\[
 \boxed{
 xO^4-(1+px)O^2+\frac{s^2}{4}=0.}
\]
If
\[
 s\neq0,
 \qquad
 \alpha^2\neq\beta^2,
\]
then both selected packets have degree four over $K(x)$.
\end{theorem}

\begin{proof}
The two equations follow by adding and subtracting the squares of
\[
 F(z)=E+zO,
 \qquad
 F(-z)=E-zO.
\]
Elimination gives the quartics.  Their quadratic discriminant is
\[
 \Delta(x)
 =
 (1-\alpha^2x)(1-\beta^2x),
\]
which is not a square under the stated hypotheses.  Hence $E^2$ generates a quadratic extension.  Its norm is $s^2x/4$, which is not a square in $K(x)$, so $E$ has degree four.  Since $EO=-s/2$, the fields generated by $E$ and $O$ coincide.
\end{proof}

The symmetric case produces an exact collision.

\begin{corollary}[Antisymmetric branch collision]
\label{cor:block10-antisymmetric-collision}
If $\beta=-\alpha\neq0$, then
\[
 \boxed{O=0,
 \qquad
 E^2=1-\alpha^2x.}
\]
For the even coordinate,
\[
 |\mathcal B_2|=4,
 \qquad
 |\mathcal C_{2,0}|=3,
 \qquad
 |\Gamma\cdot E|=2.
\]
Thus
\[
 \mathfrak c^{\rm tuple}_{2,0}=1,
 \qquad
 \mathfrak s^{\rm Gal}_{2,0}=1.
\]
For the selected odd coordinate,
\[
 |\mathcal C_{2,1}|=3,
 \qquad
 |\Gamma\cdot O|=1,
\]
so
\[
 \mathfrak c^{\rm tuple}_{2,1}=1,
 \qquad
 \mathfrak s^{\rm Gal}_{2,1}=2.
\]
\end{corollary}

\begin{proof}
Now $F(z)=F(-z)=\sqrt{1-\alpha^2z^2}$ for the selected branch, so the odd packet vanishes.  Over the split branch tuple set, the two signs are chosen independently.  The even projection has values $E,0,-E$, while the odd projection has values $0,E/z,-E/z$.  The selected Galois orbits have sizes two and one, respectively.
\end{proof}

If $\alpha=\beta$, the quadratic equation
\[
 Y^2-(1-\alpha z)^2=0
\]
is not the normalized source polynomial.  The selected branch is rational:
\[
 F_{\alpha,\alpha}(z)=1-\alpha z,
\]
so the correct source degree is one.  This is a negative control against using a nonminimal defining equation in the degree ledger.

\begin{figure}[htbp]
\centering
\includegraphics[width=.98\textwidth]{algebraic_quadratic_phase.pdf}
\caption{Catalan parity packets attain the universal quadratic bound.  Generic radicals do the same, while the antisymmetric radical separates tuple collision from selected-orbit loss.  A repeated linear factor must first be reduced to its minimal source polynomial.}
\label{fig:block10-quadratic-phase}
\end{figure}

\subsection{The directional algebraic-to-holonomic bridge}
\label{subsec:block10-algebraic-holonomic}

Algebraic functions are D-finite, but this bridge is directional unless branch data is retained.

\begin{proposition}[Algebraic derivative closure]
\label{prop:block10-algebraic-dfinite}
Let $F$ have separable minimal polynomial $P_F$ of degree $d$.  Then
\[
 F'
 =
 -\frac{\partial_zP_F(z,F)}{\partial_YP_F(z,F)}
 \in
 K(z,F).
\]
Every higher derivative lies in the same degree-$d$ function field.  Consequently, $F$ has a differential annihilator of order at most $d$.
\end{proposition}

\begin{proof}
Differentiate $P_F(z,F(z))=0$.  Separability makes $\partial_YP_F(z,F)$ nonzero in the function field.  The field is closed under the induced derivation.  The $d+1$ elements $F,F',\ldots,F^{(d)}$ are linearly dependent over $K(z)$.
\end{proof}

The annihilator can forget both the polynomial relation and the selected branch.  For example,
\[
 \sqrt{1+z}
 \quad\text{and}\quad
 -\sqrt{1+z}
\]
satisfy the same first-order differential operator.  Their initial values distinguish them.  More generally, a D-finite equation may have transcendental solutions, so an Ore chart does not reconstruct an algebraic minimal polynomial without additional recognition data.

\begin{boundarybox}
\textbf{Exact algebraic boundary.}
Algebraic atlas proves finite selected-branch normalization, an exact packet polynomial compiler, the generic $d^N$ Weil-restriction degree, a constructive coordinate eliminant, and a two-stage collision invariant.  It proves sharp packet degrees for Catalan parity and a complete collision phase for a marked quadratic-radical family.  It does not prove that every elimination output is minimal without factor selection, one canonical minimal polynomial-system dimension across arbitrary multivariate grammars, an optimal packet-degree bound using all source symmetries, or unrestricted equivalence between algebraic and D-finite presentations.
\end{boundarybox}

\subsection{Exact examples and finite verification}
\label{subsec:block10-verification}

The verifier independently constructs packet equations by substitution and reduction modulo $z^N=x$.  It checks Catalan, Motzkin, and marked radical branches at several packet widths.  It verifies the sharp Catalan quartics, the quadratic-radical trace and norm systems, the generic and antisymmetric collision phases, the rejection of nonminimal quadratic equations, and finite branch-selector certificates.

\begin{resultbox}
\textbf{Section result.}
A selected algebraic branch has a canonical minimal polynomial and a finite branch-separating jet.  Complete arithmetic packets are the solutions of an explicit $N$-equation polynomial system obtained by residue reduction.  The generic packet scheme is a finite Weil restriction of degree $d^N$, giving a constructive degree bound for every packet coordinate.  Degree loss splits exactly into tuple projection collisions and Galois stabilizer loss.  Catalan parity attains the universal quadratic bound, while the marked radical family exhibits generic sharpness, antisymmetric cancellation, and rational collapse.  The algebraic-to-holonomic bridge remains directional unless the minimal polynomial and branch selector are retained.
\end{resultbox}

\begin{table}[H]
\centering
\scriptsize
\caption{Representative exact degree and collision profiles in the Algebraic atlas algebraic packet atlas.}
\label{tab:block10-examples}
\begin{tabularx}{\textwidth}{@{}L{0.17\textwidth}L{0.14\textwidth}cccccY@{}}
\toprule
Source & Packet & $d$ & $d^N$ & degree & tuple loss & stabilizer loss & Interpretation\\
\midrule
$C=1+zC^2$ & parity $(E,O)$ & 2 & 4 & 4 & 0 & 0 & sharp Weil-restriction bound\\
$\sqrt{(1-z)(1-2z)}$ & parity even/odd & 2 & 4 & 4 & 0 & 0 & generic biquadratic packet\\
$\sqrt{1-z^2}$ & even packet & 2 & 4 & 2 & 1 & 1 & rotation alias plus descent stabilizer\\
$\sqrt{1-z^2}$ & odd packet & 2 & 4 & 1 & 1 & 2 & exact zero-address cancellation\\
$\sqrt{(1-z)^2}=1-z$ & normalized source & 1 & 1 & 1 & 0 & 0 & nonminimal quadratic equation rejected\\
\bottomrule
\end{tabularx}
\end{table}

\begin{table}[H]
\centering
\small
\caption{Exact finite validation of the Algebraic atlas algebraic packet layer.}
\label{tab:block10-validation}
\begin{tabularx}{\textwidth}{@{}L{0.34\textwidth}rY@{}}
\toprule
Validation class & Checks & Scope\\
\midrule
General Packet Compiler & 3,690 & Catalan, Motzkin, and marked quadratic radicals at packet widths two through four\\
Catalan Sharp Parity & 7,426 & exact parity system, quartic eliminants, and irreducibility ledger\\
Quadratic Radical Collision Phase & 43,900 & generic degree four, antisymmetric cancellation, and rational collapse\\
Finite Branch Selectors & 16,632 & finite jets separating explicitly colliding formal branches\\
Collision Ledger Combinatorics & 1,440 & finite tuple-collision and selected-orbit defect identities\\
\midrule
\textbf{Algebraic atlas total} & \textbf{73,088} & \textbf{0 failures}\\
Cumulative the preceding certified atlas program & 1,911,944 & frozen prior records plus recomputed Algebraic atlas\\
\bottomrule
\end{tabularx}
\end{table}


